\documentclass[10pt,a4paper]{article}

\usepackage[a4paper, margin=1.45cm]{geometry}
\usepackage[T1]{fontenc}
\usepackage{lmodern}
\usepackage{microtype}
\usepackage{booktabs}
\usepackage{tabularx}
\usepackage{array}
\usepackage{enumitem}
\usepackage{caption}
\usepackage{titlesec}
\usepackage{hyperref}
\usepackage{float}
\usepackage{graphicx}

\hypersetup{colorlinks=true, urlcolor=blue}
\captionsetup{font=footnotesize, labelfont=bf}

\titleformat{\section}{\normalsize\bfseries}{}{0em}{}[\titlerule]
\titleformat{\subsection}{\small\bfseries}{}{0em}{}
\titlespacing{\section}{0pt}{5pt}{2pt}
\titlespacing{\subsection}{0pt}{3pt}{1pt}

\setlength{\parindent}{0pt}
\setlength{\parskip}{2pt}
\setlength{\tabcolsep}{4pt}
\renewcommand{\arraystretch}{1.05}

\setlist[itemize]{leftmargin=1.2em, itemsep=0pt, topsep=1pt}

\title{\vspace{-2em}\textbf{Online Retail II --- Analysis and Report}\\[1em]
\small Carl, Cheng An, Taylor, Cyrus \quad | \quad CSP 571 --- Spring 2026 \quad | \quad April 2026}
\date{}

\begin{document}
\maketitle
\vspace{0em}

\section{Abstract}

\textit{Research Question: How can purchase timing, product associations, seasonal demand patterns, and next-day demand spike prediction be used to optimize stocking decisions and maximize revenue?}

Five modeling questions guide the analysis: (1) when do customers shop, and does this generalize internationally? (2) which product characteristics predict seasonal bulk demand? (3) if a customer bought X, what else will they buy? (4) can next-day demand spikes be predicted using temporal and demand features? (5) which customer segments exist based on purchasing behaviour, and where do re-engagement opportunities lie?

\textbf{Key findings:} UK transactions concentrate in business hours, especially Tuesday--Thursday from 08:00--17:00. Seasonal demand peaks in October--November, with invoice frequency as the dominant stocking signal. Apriori produced 1,268 co-purchase rules mapping clearly to product families. For next-day spike prediction, GBM at a decision threshold of 0.3 achieved the highest UK recall at 65.5\%. Customer segmentation identified four distinct RFM-based tiers, with 1,230 dormant high-value customers representing a re-engagement opportunity exceeding \pounds250,000. All models degraded on non-UK data, suggesting region-specific recalibration is required.

\section{Data Processing}

The raw Online Retail II data exhibited several quality issues, including missing product descriptions, inconsistent product labeling, and non-retail transaction records such as administrative entries, returns, financial adjustments, and free-item transactions. The data processing pipeline was designed to isolate high-confidence retail sales while preserving non-standard records in separate datasets for audit and validation.

\subsection{Transaction Filtering}

The \texttt{InvoiceNo}, \texttt{quantity}, and \texttt{price} variables were used to classify transaction types. Invoices beginning with \texttt{C} were treated as cancellations or returns, while invoices beginning with \texttt{A} were treated as administrative/accounting records. Rows with negative quantity and zero price were classified as inventory adjustments, and rows with positive quantity and zero price were classified as free or promotional items. Standard clean sales were defined as records with positive quantity and positive price, excluding cancellations, administrative rows, adjustments, free items, shipping fees, gift vouchers, DCGS entries, special operational codes, and manual rows. Manual entries were isolated as product-like sales, allowing them to be optionally included or excluded depending on the analysis.

\subsection{Product Description and Stock Code Standardization}

The \texttt{stock\_code} variable was treated as the primary product identifier. This was necessary because product descriptions were not always consistent: identical products sometimes appeared with spelling differences, spacing differences, or operational labels such as ``found'', ``lost'', or ``damaged''. The pipeline first filled missing descriptions when a stock code had exactly one known description. It then identified noise labels and assigned each product a canonical description using the most frequent non-noise description associated with that \texttt{stock\_code}. This approach assumes that each stock code corresponds to one underlying product, while acknowledging that special stock codes such as shipping charges, gift vouchers, and administrative codes require separate handling.

\subsection{Data Quality Improvements}

The cleaning process produced measurable improvements in dataset quality. The raw data contained over 4{,}000 missing product descriptions, which were reduced to 363 remaining missing descriptions after cleaning and standardization. In addition, approximately 6\% of all rows required description standardization, indicating moderate inconsistency in product labeling. After rule-based filtering, approximately 97\% of the full dataset was classified as high-confidence retail sales. The complete cleaned dataset was stored as \texttt{cleaned\_retail\_all}, while derived tables such as \texttt{clean\_product\_sales}, \texttt{manual\_product\_sales}, \texttt{cancellations}, and \texttt{inventory\_adjustments} were stored separately for downstream analysis.

\section{Data Analysis}
\subsection{Summary Statistics}

Summary statistics are computed on the clean-sales subset (positive quantity and price, excluding cancellations and non-product codes).

\begin{table}[H]
\centering
\caption{Key summary statistics (clean sales)}
\footnotesize
\begin{tabular}{lr}
\toprule
Metric & Value \\
\midrule
Rows (clean sales) & $\sim$1.04M \\
Unique invoices & $\sim$40K \\
Unique customers & $\sim$5.9K \\
Unique products & $\sim$4.9K \\
Total revenue & $\sim$\pounds{}21M \\
Avg revenue / invoice & $\sim$\pounds{}523 \\
\bottomrule
\end{tabular}
\end{table}

\subsection{Data Visualization}

\begin{figure}[H]
    \centering
    \includegraphics[width=0.5\textwidth]{figures/transaction_value_distribution.png}
    \caption{Distribution of transaction-level revenue. The distribution is highly right-skewed, indicating that most purchases are low-value while a smaller number of bulk transactions generate disproportionately large revenue.}
    \label{fig:transaction_value_distribution}
\end{figure}

\begin{figure}[H]
    \centering
    \includegraphics[width=0.5\textwidth]{figures/monthly_sales_trend.png}
    \caption{Monthly sales trend for clean retail transactions. Revenue shows strong temporal variation, with clear seasonal peaks that motivate the use of time-based features in downstream modeling.}
    \label{fig:monthly_sales_trend}
\end{figure}

\begin{figure}[H]
    \centering
    \includegraphics[width=0.5\textwidth]{figures/top_products_revenue.png}
    \caption{Top 10 products by total revenue. Revenue is concentrated among a small set of products, suggesting that product-level demand history is an important signal for sales analysis and recommendation tasks.}
    \label{fig:top_products_revenue}
\end{figure}

The visualizations reveal three important patterns. First, transaction revenue is highly right-skewed, meaning that most purchases are small while a limited number of bulk transactions contribute disproportionately to total revenue. Second, monthly sales show clear seasonal variation, suggesting that calendar-based features are important for demand modeling. Third, revenue is concentrated among a small number of products, supporting the use of product-level features such as historical demand, invoice frequency, and revenue contribution.

\subsection{Key Insights}

Transaction-level revenue is highly right-skewed, with most purchases being low-value and a small number of bulk transactions contributing disproportionately to total revenue. This imbalance introduces sensitivity to outliers, suggesting the need for modeling approaches that can handle skewed target distributions.

Sales exhibit clear seasonal patterns, with a pronounced increase in activity during October and November. This highlights the importance of temporal features such as month and week-of-year for demand modeling and forecasting tasks.

Demand is highly concentrated among a small subset of products, indicating that a limited number of items drive a significant share of overall sales. This concentration suggests that product-level features, including historical demand and purchase frequency, are likely to be strong predictors.

The dataset is heavily dominated by UK transactions, with international markets representing a smaller and more variable share of activity. This imbalance motivates the use of out-of-distribution testing to evaluate model generalization across regions.


Key observations:
\begin{itemize}
\item Sales are highly skewed, with a small number of products driving a large share of revenue.
\item The dataset contains a large number of low-value transactions and fewer high-value bulk purchases.
\item Returns (C-invoices) account for a non-trivial portion of the dataset, confirming the importance of filtering.
\item Missing CustomerID values remain a limitation, with approximately 20--25\% of transactions lacking customer identifiers, restricting customer-level analysis.
\end{itemize}

\subsection{Feature Extraction}

Features were derived to support downstream modeling tasks, including customer segmentation and demand prediction.

Customer-level features follow the RFM framework:
\begin{itemize}
\item Recency: days since last purchase
\item Frequency: number of invoices
\item Monetary: total spend
\end{itemize}

Transaction-level features include:
\begin{itemize}
\item TotalPrice (Quantity $\times$ UnitPrice)
\item Month, DayOfWeek, Hour
\item Country indicators (e.g., UK vs non-UK)
\end{itemize}

Product-level features capture demand concentration and purchasing behavior, including total quantity sold, number of customers, and revenue contribution. These features directly support the modeling tasks described in the subsequent sections.


\section{Model Training}

\subsection{Train/Test Strategy and Feature Engineering}
The UK market represents approximately 92\% of invoice volume and is used as the primary training domain. An 80/20 split creates in-distribution UK train/test sets, while non-UK transactions are held out as an out-of-distribution (OOD) test set.

For peak-hour classification, \texttt{IsPeakHour} flags invoices above the 67th percentile of hourly average revenue, computed using UK training data only. Predictors include \texttt{TotalRevenue}, \texttt{TotalQty}, \texttt{NumItems}, \texttt{DayOfWeek}, and \texttt{MonthNum}. For seasonal bulk demand, records are aggregated to product--month level, with \texttt{TotalQty} predicted from \texttt{MonthNum}, \texttt{AvgPrice}, and \texttt{NumInvoices}. Apriori uses product-description sets per UK invoice directly.

For demand spike prediction, transaction data is aggregated daily. The target \texttt{next\_day\_spike} identifies whether the following day's quantity is in the top 25\%, corresponding to approximately 18,460 units. Predictors include lag features, 7-day rolling averages of orders and quantity, calendar variables including week of year, day of week, and their interaction, and a weekend indicator. All predictors use only information available before the prediction day to avoid leakage.

For customer segmentation, RFM features --- Recency, Frequency, and Monetary value --- are constructed at the customer level using 5,851 customers after cleaning. Due to right skew, a log transformation is applied before PCA-based dimensionality reduction. The first three principal components explain 91.6\% of variance and are used as inputs to both K-Means and Hierarchical Clustering.

\subsection{Model Selection}
\begin{table}[H]
\centering
\caption{Model selection and rationale}
\footnotesize
\begin{tabularx}{\textwidth}{
>{\raggedright\arraybackslash\bfseries}p{2.7cm}
>{\raggedright\arraybackslash}p{3.4cm}
>{\raggedright\arraybackslash}X}
\toprule
Question & Models & Rationale \\
\midrule
When do customers shop?
& Logistic Regression, LDA
& Binary classification; LDA checks whether a linear decision boundary is sufficient \\
Seasonal bulk buying
& GAM, Regression Spline, Random Forest
& Captures nonlinear seasonality; RF provides variable importance \\
If bought X, buy Y
& Apriori association rules
& Models co-purchase probability using support, confidence, and lift \\
Demand spike prediction
& Decision Tree, Random Forest, GBM
& Tree-based models capture nonlinear interactions; GBM at threshold 0.3 prioritizes recall under cost-sensitive evaluation, catching 65.5\% of true spike days \\
Customer segmentation
& K-Means, Hierarchical Clustering
& RFM-based clustering identifies actionable customer tiers; both methods independently confirm four segments \\
\bottomrule
\end{tabularx}
\end{table}

\section{Model Validation and Performance}
\textbf{When do customers shop?} Logistic Regression and LDA both achieve 32.3\% UK error and approximately 41.4\% non-UK error. Their near-identical performance suggests a mostly linear boundary, while OOD degradation shows that shopping-hour patterns do not fully transfer internationally.
"When do customers shop?" has a pattern that is consistent with a wholesale/trade customer base rather than a consumer retailer.

\begin{figure} [H]
    \centering
    \includegraphics[width=0.5\linewidth]{when_do_customers_shop.png}
    \caption{The bulk of transactions happen between roughly 8am and 5pm on weekdays, which makes sense for a B2B-leaning online retailer. Weekends are almost completely dark, meaning virtually no trading activity on Saturdays and Sundays.}
    \label{fig:when_do_customers_shop}
\end{figure}

\textbf{Seasonal products and bulk buying.} Cross-validation selected \texttt{df = 11} for the regression spline, indicating a flexible seasonal pattern. GAM ANOVA confirms all predictors are significant ($p < 0.001$). Random Forest importance ranks \texttt{NumInvoices} $\gg$ \texttt{AvgPrice} $>$ \texttt{MonthNum}, meaning reorder frequency is stronger than price or month alone. RF test MSE rises from 138,431 on UK data to 293,590 on non-UK data.

\textbf{If bought X, buy Y.} Apriori on UK invoices produced 1,268 rules using minimum support 0.01 and minimum confidence 0.20. Median lift is 8.2 and maximum lift is 47.7. The strongest rules map to product families such as Poppy's Playhouse sets and coordinated tableware. Most rules transfer internationally with small confidence drops, but some pairs show UK-specific behavior. The exported rules CSV powers the interactive \texttt{Next\_likely\_product.py} recommendation tool.

\textbf{Demand spike prediction.} Three classification models were trained using time series features including lag variables, 7-day rolling averages, and calendar features. A demand spike is defined as any day in the top 25\% of daily sales, approximately above 18,460 units. Random Forest achieved the highest accuracy at 85.7\%, but GBM with a lowered decision threshold of 0.3 was selected as the recommended model under a cost-sensitive objective, where missing a spike is more costly than a false positive. At this threshold, GBM captures 65.5\% of true spike days. The strongest predictors are week of year (27\%), followed by 7-day rolling averages of orders and quantity, indicating that recent sales momentum carries predictive power independently of seasonality. The interaction between week of year and day of week adds further signal, while weekend status has negligible impact.

\begin{table}[H]
\centering
\caption{Demand spike model performance}
\footnotesize
\begin{tabular}{lcccccc}
\toprule
& \multicolumn{3}{c}{UK Test} & \multicolumn{3}{c}{Non-UK OOD} \\
\cmidrule(lr){2-4}\cmidrule(lr){5-7}
Model & Acc. & Prec. & Rec. & Acc. & Prec. & Rec. \\
\midrule
Decision Tree & 82.4 & 78.6 & 37.9 & 71.4 & 33.9 & 14.0 \\
Random Forest & 85.7 & 75.0 & 62.1 & 73.7 & 31.2 & 3.5 \\
\textbf{GBM} & \textbf{84.9} & \textbf{70.4} & \textbf{65.5} & \textbf{70.1} & \textbf{37.7} & \textbf{28.0} \\
\bottomrule
\end{tabular}
\end{table}

\textbf{Customer segmentation.} RFM profiles were constructed for 5,334 customers. Recency measures how recently a customer last purchased, Frequency measures how often they order, and Monetary value measures total spend. Both K-Means and Hierarchical Clustering independently converged on four segments. Segment 1 comprises 1,566 loyal high-value customers averaging 16 orders per year and nearly \pounds9,000 in total spend. Segment 3 is the largest group with 1,917 active mid-tier customers averaging 4 orders per year. Segment 4 contains 1,138 lapsed low-value customers who purchased once or twice and did not return. Segment 2 is the most strategically notable: 1,230 customers with an average order value of \pounds567, comparable to the best customers, but who have placed only 2 orders and have been inactive for over 400 days.

\begin{table}[H]
\centering
\caption{Customer Segment Profiles (K-Means, K = 4)}
\footnotesize
\begin{tabular}{lccccc}
\toprule
Segment & Customers & Avg.\ Recency (days) & Avg.\ Frequency & Avg.\ Monetary (\pounds) & Avg.\ Order Value (\pounds) \\
\midrule
1 — Loyal High-Value    & 1,566 & 55  & 15.9 & 8,898 & 546 \\
2 — Dormant High-Value  & 1,230 & 400 & 1.9  & 1,045 & 567 \\
3 — Active Mid-Tier     & 1,917 & 90  & 4.0  & 948   & 261 \\
4 — Lapsed Low-Value    & 1,138 & 372 & 1.6  & 202   & 133 \\
\bottomrule
\end{tabular}
\end{table}

\section{Conclusion}

UK stocking decisions can be improved by aligning inventory with October-November demand peaks, stocking high-lift co-purchase families together, and using GBM-based spike alerts for proactive restocking. Across models, invoice frequency and rolling demand momentum are more reliable signals than price, weekends, or single-day activity. Customer segmentation further enables targeted re-engagement: Segment 2 alone represents a recoverable revenue opportunity exceeding \pounds250,000 if 20\% of dormant high-value customers are reactivated.

\textbf{Positive results:}
\begin{itemize}
  \item 1,067,371 records were cleaned and modeled with reproducible workflows.
  \item UK shopping concentrates in business hours with negligible weekend activity.
  \item Apriori extracted 1,268 interpretable co-purchase rules.
  \item GBM (threshold 0.3) achieved 65.5\% recall for UK demand spike prediction, capturing the majority of high-demand days before they occur.
  \item RFM clustering identified four actionable customer segments across 5,334 customers, with both K-Means and Hierarchical Clustering independently agreeing on the same structure.
  \item OOD testing consistently revealed international transfer degradation.
\end{itemize}

\textbf{Caveats and recommendations:}
\begin{itemize}
  \item Peak-hour error remains high, so basket features alone are weak timing predictors.
  \item 22.8\% missing \texttt{CustomerID} values limit customer-level segmentation.
  \item Demand spike models do not generalize reliably to non-UK markets.
  \item Niche products below Apriori support thresholds are excluded from rules.
  \item A sequential rule mining approach (like SPADE or PrefixSpan) would tell what customers buy on their next visit after buying X. Knowing this could help with retention.
  \item Region-specific models should replace a single global model.
  \item A deep learning model for predicting upcoming orders and revenue would need more data and richer features, but it's a natural extension of what we've built here.
\end{itemize}

\section{Data Sources and Code}

\footnotesize
\textbf{Dataset:} Online Retail II Dataset, UCI Machine Learning Repository. 
\url{https://archive.ics.uci.edu/dataset/502/online+retail+ii}

\textbf{Code:} \url{https://github.com/CAWang16/business_data_prep/tree/main}

\section{Bibliography}

\footnotesize
Bult, J.R., and Wansbeek, T. (1995). Optimal selection for direct mail. \textit{Marketing Science}, 14(4), 378--394.

Chen, D., Sain, S.L., and Guo, K. (2012). Data mining for the online retail industry. \textit{Journal of Database Marketing \& Customer Strategy Management}, 19(3), 197--208.

Rahm, E., and Do, H.H. (2000). Data cleaning: Problems and current approaches. \textit{IEEE Data Engineering Bulletin}, 23(4), 3--13.

\end{document}
